This paper presented RADAR, a hybrid neural architecture for robust AI-generated text
detection that combines deep contextual semantics with handcrafted stylometric signals.
Built on ModernBERT-large, RADAR couples a frozen or fine-tuned encoder (ARSE) with a
stylometric feature extractor (SLPE) and a cross-attention fusion module that lets
semantic and stylometric representations interact in a principled way before a single
binary classifier produces a calibrated confidence score.

We trained RADAR through a two-stage curriculum on a heterogeneous dataset combining the
official PAN-2026 training data and a balanced subsample of the RAID benchmark, which
encompasses a wide variety of AI generators and text-level adversarial perturbations.
The final end-to-end model achieves Mean~=~0.997 on the PAN-2026 validation set,
surpassing all published baselines, and Mean~=~0.972 on the RAID test set, demonstrating
strong cross-domain generalisation.

The ablation experiments reveal several key findings. First, multi-source training is
indispensable for robust performance: models trained on a single source degrade
substantially under domain shift, with Exp.~3 (PAN-2026 only) dropping from
Mean~=~0.985 in-domain to Mean~=~0.681 on RAID. Second, the two-stage curriculum
consistently outperforms single-stage training: unfreezing the encoder only after the
non-encoder modules have converged leads to stable and fast convergence of the full
model. Third, stylometric features complement semantic representations in a way that
improves robustness to adversarial obfuscation attacks such as synonym substitution
and homoglyph injection, which systematically fool purely lexical models.

\paragraph{Limitations.}
Although RADAR achieves strong performance on the evaluated benchmarks, several
limitations warrant caution in deployment. The model is trained predominantly on
English texts; its effectiveness on other languages or mixed-language content is unknown.
The architecture's maximum input length of 512 tokens means that the signal from longer
documents is truncated. Furthermore, while RAID covers a diverse range of adversarial
attacks, new obfuscation strategies may emerge that are not yet represented in the training
distribution. Finally, the stylometric features are extracted via deterministic rules;
learned stylometric representations~\cite{halvani:2017} could potentially offer richer
signals.

\paragraph{Future Work.}
Promising directions for future work include: (i) extending the evaluation to multilingual
texts and multilingual base encoders; (ii) incorporating longer-range context through
hierarchical encoding or sparse attention; (iii) exploring adversarial training with
additional obfuscation types at training time; and (iv) applying uncertainty calibration
techniques to further improve the reliability of confidence scores for human-in-the-loop
review workflows.

In summary, RADAR demonstrates that combining pre-trained language model representations
with lightweight, theoretically motivated stylometric signals through a principled
cross-attention fusion mechanism is a highly effective and generalisable approach to
AI-generated text detection in adversarial conditions.
